\BourbakiExercisesSection

\begin{enumerate}
\def\labelenumi{\arabic{enumi}.}
\item
  Determine all the ring structures on a set of \(n\) elements for \(2 \leqslant n \leqslant 7\) and also the ideals of these rings.
\item
  Let A be a ring.
\end{enumerate}

\begin{enumerate}
\def\labelenumi{(\alph{enumi})}
\item
  The mapping which associates with \(a \in A\) the left homothety \(x \mapsto ax\) is an isomorphism of A onto the ring of endomorphisms of the additive group of A which commute with right homotheties.
\item
  Let M be the set of sequences \((a, b, c, d)\) of four elements of A which may be represented in the form of a table or matrix \(\begin{pmatrix} a & b \\ c & d \end{pmatrix}\) . With every element \(\begin{pmatrix} a & b \\ c & d \end{pmatrix}\) of M we associate the mapping \((x, y) \mapsto (ax + by, cx + dy)\) of \(\mathbf{A} \times \mathbf{A}\)
\end{enumerate}

into itself. Show that a bijection is thus obtained of M onto the ring B of endomorphisms of the commutative group A × A which commute with the operations \((x, y) \mapsto (xr, yr)\) for all \(r \in A\) . Henceforth M is identified with the ring B. Calculate the product of two matrices.

\begin{enumerate}
\def\labelenumi{(\alph{enumi})}
\setcounter{enumi}{2}
\tightlist
\item
  Show that the subset \(\mathbf{T}\) of \(\mathbf{M}\) consisting of the matrices of the form \(\begin{pmatrix} a & b \\ 0 & c \end{pmatrix}\) is the subring of \(\mathbf{M}\) which leaves stable the subgroup \(\mathbf{A} \times 0\) of \(\mathbf{A} \times \mathbf{A}\) .
\end{enumerate}

The mappings \(\begin{pmatrix} a & b \\ 0 & c \end{pmatrix} \mapsto a\) and \(\begin{pmatrix} a & b \\ 0 & c \end{pmatrix} \mapsto c\) are homomorphisms of T into A; let \(a\) and \(b\) be their kernels. Show that \(a + b = T\) , that \(ab = \{0\}\) and that \(ba = b \cap a\) is \(\neq \{0\}\) if \(A \neq \{0\}\) (cf.~no. 9, Proposition 6).

\begin{enumerate}
\def\labelenumi{(\alph{enumi})}
\setcounter{enumi}{3}
\tightlist
\item
  If \(\mathbf{A} = \mathbf{Z} / 2\mathbf{Z}\) , \(\mathbf{T}\) is not commutative, but the multiplicative group of invertible elements of \(\mathbf{T}\) is commutative.
\end{enumerate}

\begin{enumerate}
\def\labelenumi{\arabic{enumi}.}
\setcounter{enumi}{2}
\item
  Let A be a ring and \(a \in A\) . If there exists one and only one \(a' \in A\) such that \(aa' = 1\) , a is invertible and \(a' = a^{-1}\) (show first that a is left cancellable, then consider the product \(aa'a\) ).
\item
  Let A be a pseudo-ring in which every additive subgroup of A is a left ideal of A.
\end{enumerate}

\begin{enumerate}
\def\labelenumi{(\alph{enumi})}
\item
  If \(\mathbf{A}\) is a ring, \(\mathbf{A}\) is canonically isomorphic to \(\mathbf{Z} / n\mathbf{Z}\) for some suitable integer \(n\) .
\item
  If A has no unit element and every non-zero element is cancellable, A is isomorphic to a pseudo-subring of Z.
\end{enumerate}

\begin{enumerate}
\def\labelenumi{\arabic{enumi}.}
\setcounter{enumi}{4}
\tightlist
\item
  Let \(\mathbf{M}\) be a monoid, \(\mathbf{Z}^{(\mathbf{M})}\) the free commutative group with basis \(\mathbf{M}\) and \(u\colon \mathbf{M}\to \mathbf{Z}^{(\mathbf{M})}\) the canonical mapping.
\end{enumerate}

\begin{enumerate}
\def\labelenumi{(\alph{enumi})}
\item
  Show that there exists a unique multiplication on \(\mathbf{Z}^{(\mathbf{M})}\) such that \(\mathbf{Z}^{(\mathbf{M})}\) is a ring and such that \(u\) is a monoid morphism of the monoid \(\mathbf{M}\) into the multiplicative monoid of the ring \(\mathbf{Z}^{(\mathbf{M})}\) .
\item
  Let A be a ring and \(v: \mathbf{M} \to \mathbf{A}\) a monoid morphism of the monoid M into the multiplicative monoid of A. There exists a unique ring morphism \(f: \mathbf{Z}^{(\mathbf{M})} \to \mathbf{A}\) such that
\end{enumerate}

\[
f \circ u = v.
\]

\begin{enumerate}
\def\labelenumi{\arabic{enumi}.}
\setcounter{enumi}{5}
\tightlist
\item
  Let A be a commutative ring and M a submonoid of the additive group of A. Let n be an integer \textgreater0 such that \(m^{n} = 0\) for \(m \in M\) .
\end{enumerate}

\begin{enumerate}
\def\labelenumi{(\alph{enumi})}
\item
  If \(n!\) is not a divisor of zero in \(\mathbf{A}\) , the product of a family of \(n\) elements of \(\mathbf{M}\) is zero.
\item
  Show that the conclusion in (a) is not necessarily true if \(n!\) is a divisor of zero in A.
\end{enumerate}

\begin{enumerate}
\def\labelenumi{\arabic{enumi}.}
\setcounter{enumi}{6}
\tightlist
\item
  For all \(n \in \mathbf{N}^*\) ,
\end{enumerate}

\[
n! = \sum_ {\nu = 0} ^ {n - 1} (- 1) ^ {\nu} \binom {n} {\nu} (n - \nu) ^ {n}
\]

(use no. 2, Proposition 2).

\begin{enumerate}
\def\labelenumi{\arabic{enumi}.}
\setcounter{enumi}{7}
\item
  In a ring, the right ideal generated by a left ideal is a two-sided ideal.
\item
  In a ring, the right annihilator of a right ideal is a two-sided ideal.
\item
  In a ring A, the two-sided ideal generated by the elements xy - yx, where x and y run through A, is the smallest of the two-sided ideals a such that A/a is commutative.
\end{enumerate}

¶ 11. In a ring A, a two-sided ideal a is said to be irreducible if there exists no ordered pair of two-sided ideals b, c distinct from a and such that a = b ∩ c.

\begin{enumerate}
\def\labelenumi{(\alph{enumi})}
\item
  Show that the intersection of all the irreducible two-sided ideals of A reduces to 0 (note that the set of two-sided ideals containing no element \(a \neq 0\) is inductive and apply Zorn's Lemma).
\item
  Deduce that every two-sided ideal \(\mathfrak{a}\) of \(\mathbf{A}\) is the intersection of all the irreducible ideals which contain it.
\end{enumerate}

\begin{enumerate}
\def\labelenumi{\arabic{enumi}.}
\setcounter{enumi}{11}
\tightlist
\item
  Let A be a ring and \((\mathfrak{a}_{i})_{i\in I}\) a finite family of left ideals of A such that the additive group of A is the direct sum of the additive groups \(a_{i}\) . If \(1=\sum_{i\in I}e_{i}(e_{i}\in\mathfrak{a}_{i})\) , then \(e_{i}^{2}=e_{i}, e_{i}e_{j}=0\) if \(i\neq j\) and \(a_{i}=Ae_{i}\) (write x=x.1 for all \(x\in A\) ). Conversely, if \((e_{i})_{i\in I}\) is a finite family of idempotents of A such that \(e_{i}e_{j}=0\) for \(i\neq j\) and \(1=\sum_{i}e_{i}\) , then A is the direct sum of the left ideals \(Ae_{i}\) . For the ideals \(Ae_{i}\) to be two-sided, it is necessary and sufficient that the \(e_{i}\) be in the centre of A.
\end{enumerate}

¶ 13. Let A be a ring and e an idempotent of A.

\begin{enumerate}
\def\labelenumi{(\alph{enumi})}
\item
  Show that the additive group of \(\mathbf{A}\) is the direct sum of the left ideal \(\mathfrak{a} = \mathbf{A}e\) and the left annihilator \(\mathfrak{b}\) of \(e\) (note that, for all \(x\in \mathbf{A}\) , \(x - xe\in \mathfrak{b}\) ).
\item
  Every right ideal \(\mathfrak{d}\) of \(\mathbf{A}\) is the direct sum of \(\mathfrak{d} \cap \mathfrak{a}\) and \(\mathfrak{d} \cap \mathfrak{b}\) .
\item
  If \(\mathbf{Ae} = e\mathbf{A}\) , \(\mathfrak{a}\) and \(\mathfrak{b}\) are two-sided ideals of \(\mathbf{A}\) and define a direct decomposition of \(\mathbf{A}\) .
\end{enumerate}

\begin{enumerate}
\def\labelenumi{\arabic{enumi}.}
\setcounter{enumi}{13}
\item
  Let A be a commutative ring in which there is only a finite number n of divisors of zero. Show that A has at most \((n+1)^{2}\) elements. (Let \(a_{1},\ldots,a_{n}\) be the divisors of zero. Show that the annihilator \(J_{1}\) of \(a_{1}\) has at most \(n+1\) elements and that the quotient ring A/ \(J_{1}\) has at most \(n+1\) elements.)
\item
  A ringoid is a set E with two laws of composition: (a) an associative multiplication \(xy\) ; (b) a law written additively, not everywhere defined (§ 1, no. 1) and satisfying the following conditions:
\end{enumerate}

\begin{enumerate}
\def\labelenumi{(\arabic{enumi})}
\item
  it is commutative (in other words, if \(x + y\) is defined, so is \(y + x\) and \(x + y = y + x\) ; \(x\) and \(y\) are then said to be addible);
\item
  if \(x\) and \(y\) are addible, for \(x + y\) and \(z\) to be addible it suffices that \(x\) and \(y\) on the one hand, and \(y\) and \(z\) on the other, be addible; then \(x\) and \(y + z\) are addible and \((x + y) + z = x + (y + z)\) ;
\item
  there exists an identity element 0;
\item
  if \(x\) and \(z\) on the one hand, \(y\) and \(z\) on the other, are addible and \(x + z = y + z\) , then \(x = y\) ;
\item
  if \(x, y\) are addible, so are \(xz\) and \(yz\) (resp. \(zx\) and \(zy\) ) and
\end{enumerate}

\[
(x + y) z = x z + y z
\]

(resp. \(z(x + y) = zx + zy)\) for all \(z \in \mathbf{E}\) .

Every ring is a ringoid.

\begin{enumerate}
\def\labelenumi{(\alph{enumi})}
\item
  Examine how the definitions and results of §8 and the above exercises extend to ringoids (a left ideal of a ringoid E is a subset \(\mathfrak{a}\) of E which is stable under addition and such that \(\mathbf{E}.\mathfrak{a}\subset \mathfrak{a}\) ).
\item
  Let G be a group with operators and f and g two endomorphisms of G. For the mapping \(x \mapsto f(x)g(x)\) to be an endomorphism of G, it is necessary and sufficient that every element of the subgroup \(f(\mathrm{G})\) be permutable with every element of \(g(\mathrm{G})\) ; if this endomorphism is then denoted by \(f + g\) and fg is the composite endomorphism \(x \mapsto f(g(x))\) , show that the set E of endomorphisms of G with these laws of composition is a ringoid with unit element; for E to be a ring, it is necessary and sufficient that G be commutative.
\end{enumerate}

For an element \(f \in E\) to be addible to all the elements of E, it is necessary and sufficient that \(f(G)\) be contained in the centre of G; the set N of these endomorphisms is a pseudo-ring called the kernel of the ringoid E.

An endomorphism \(f\) of \(G\) is called normal if it is permutable with all the inner automorphisms of \(G\) ; for every normal stable subgroup \(H\) of \(G, f(G)\) is then a normal stable subgroup of \(G\) . Show that the set \(D\) of normal endomorphisms of \(G\) forms a subringoid of \(E\) and that the kernel \(N\) is a two-sided ideal in \(D.†\)

\begin{enumerate}
\def\labelenumi{\arabic{enumi}.}
\setcounter{enumi}{15}
\tightlist
\item
  Give examples of ideals \(a, b, c\) in the ring \(\mathbf{Z}\) such that \(ab \neq a \cap b\) and \((a + b)(a + c) \neq a + bc\) .
\end{enumerate}

